\section{Term Type Generation}
\label{sec:method}

We build the taxonomy in stages. The first stage, which this section describes,
asks a single question of the sources. Along which dimensions do they classify
crypto-assets, the systems that record them, the actors that handle them, the
activities those actors perform, and the obligations attached to each? We call
these dimensions \emph{term types}. Later stages place terms under the term
types and link the term types to one another (\cref{sec:method:hierarchy}).
A term type exists in our output only if the model can quote a passage that
uses it as a criterion, and we check every quote against the corpus after the
run.

\begin{figure*}[t]
  \centering
  \todo{Overview figure: PDFs $\rightarrow$ definitions-first extraction
  $\rightarrow$ XML corpus behind a prompt cache $\rightarrow$ one extraction
  call $\rightarrow$ JSON term types $\rightarrow$ citation verifier
  $\rightarrow$ verified catalogue.}
  %\includegraphics[width=\textwidth]{overview.pdf}
  \caption{Overview of the term type generation stage of \sys. \todo{caption}}
  \label{fig:overview}
\end{figure*}

\subsection{A Definitions-First Corpus}
\label{sec:method:corpus}

We do not give the model whole statutes. A regulation such as \src{MiCA}
defines its vocabulary in a few articles and spends the rest of its length on
obligations that refer back to those definitions. The classification criteria
we are after live in the definitions articles, the glossaries, the
classification conditions and the descriptions of business models. We take
those sections and nothing else. \Cref{tab:corpus} lists the result: 21 legal
instruments from the EU, the UK, the US and six international bodies, plus two
academic taxonomies of blockchains~\cite{tasca2019taxonomy} and of
crypto-assets~\cite{zhang2026cryptotaxonomy} and one paper describing a
domain corpus~\cite{hernandez2026dltcorpus}, with reference lists and large
tables removed.

\begin{table}[t]
  \centering
  \caption{\corpus: the 24 source documents, the sections we kept from each,
  their length, and the number of Task A term types that cite them.}
  \label{tab:corpus}
  \scriptsize
  %% Generated by paper/scripts/task_a_tables.py -- do not edit by hand.
\setlength{\tabcolsep}{3.5pt}
\begin{tabular}{@{}rllrrr@{}}
    \toprule
    \# & Jur. & Source & Sect. & Words\,(k) & Types \\
    \midrule
    1 & Acad. & Tasca \& Tessone taxonomy & 2 & 16.5 & 21 \\
    2 & Acad. & Zhang et al.\ taxonomy & 2 & 9.1 & 24 \\
    3 & Acad. & DLT-Corpus paper & 1 & 6.2 & 1 \\
    4 & EU & MiCA, Reg.\ 2023/1114 & 1 & 2.7 & 25 \\
    5 & EU & AMLR, Reg.\ 2024/1624 & 1 & 4.4 & 11 \\
    6 & EU & Transfer of Funds, Reg.\ 2023/1113 & 1 & 1.0 & 11 \\
    7 & EU & DORA, Reg.\ 2022/2554 & 1 & 2.3 & 5 \\
    8 & UK & FSMA (Cryptoassets) Regs 2026 & 2 & 4.6 & 24 \\
    9 & UK & CARF Regs 2025 & 1 & 0.6 & 3 \\
    10 & US & CLARITY Act (Senate-reported) & 2 & 1.4 & 13 \\
    11 & US & CLARITY Act (House-passed) & 1 & 4.8 & 20 \\
    12 & US & FIT21, H.R.\ 4763 & 1 & 4.7 & 20 \\
    13 & US & GENIUS Act, Pub.\ L.\ 119-27 & 1 & 1.7 & 20 \\
    14 & US & Executive Order 14178 & 1 & 0.1 & 3 \\
    15 & US & PWG Digital Asset Markets Report & 1 & 12.2 & 22 \\
    16 & US & FinCEN CVC Guidance 2019 & 3 & 7.2 & 17 \\
    17 & Intl & FATF VA/VASP Guidance 2021 & 2 & 9.1 & 27 \\
    18 & Intl & FATF Targeted Update 2025 & 1 & 1.7 & 10 \\
    19 & Intl & FATF Targeted Update 2026 & 1 & 2.6 & 12 \\
    20 & Intl & Basel SCO60 & 2 & 3.1 & 15 \\
    21 & Intl & FSB Recommendations 2023 & 1 & 0.5 & 14 \\
    22 & Intl & IOSCO Recommendations 2023 & 1 & 0.7 & 11 \\
    23 & Intl & UN Cybercrime Convention & 1 & 0.7 & 7 \\
    24 & Intl & Budapest Convention & 1 & 0.2 & 3 \\
    \midrule
    \multicolumn{3}{@{}l}{Total} & 32 & 97.8 & 93 \\
    \bottomrule
\end{tabular}

\end{table}

Each PDF is converted with \texttt{pdftotext} and cut by line ranges or by
regular-expression markers that we specify once per document. The US bills
need the layout-preserving mode because their two-column line numbering
scrambles the plain output, and we drop page footers such as the bill number
that would otherwise leak into every section. The Senate-reported print of the
\src{CLARITY Act} reproduces the House-passed text before the Senate
substitute, so we take Title~I once, from the House file, and only the
substitute's definitions and its section on network tokens from the Senate
file. The corpus is one XML-tagged text file. Every document carries an integer
index and metadata (type, jurisdiction, title, year) and is split into
\texttt{<section label="...">} blocks, so a citation can name a document and a
section unambiguously. The file holds 32 sections and 97.8 thousand words,
which the model's tokeniser counts as 235,298 tokens.

\subsection{The Extraction Task}
\label{sec:method:task}

The system prompt defines the objects we want. A \emph{term} is a word or
phrase the corpus defines or classifies, such as \term{crypto-asset} or
\term{permissioned ledger}. A \emph{term type} is a question along which the
corpus's definitions discriminate between things, such as \ttype{reference
asset} or \ttype{consensus mechanism}, and it has \emph{values}, the categories
a thing can take on that dimension. A term type is \emph{grounded} when at
least one passage uses it as a criterion or defines several terms that differ
along it. The model is told that the corpus is its only source of truth and
that outside knowledge may be used to recognise what the corpus says, never to
add to it.

Three grounding rules follow. Every term type and every value must carry at
least one citation of the form \{document index, section label, quote\}; the
quote must be verbatim and at most thirty words; a candidate that cannot be
grounded must be omitted. Four quality rules shape the output: one question per
term type, with differently named versions of the same question merged and the
other names kept as aliases; jurisdictional disagreement about values or
thresholds recorded in a notes field with citations for each side; preference
for dimensions that recur across sources or that a source uses to draw a legal,
prudential or technical boundary, while still keeping an unusual dimension that
a single source grounds clearly; and JSON only.

The user prompt states the scope, which covers assets, DLT systems, actors,
activities, and obligations or risks when the sources themselves use them as
criteria, and gives the JSON schema. Each term type has an identifier, a name,
the category it applies to, the question it answers, a short description, its
aliases, its values with their own definitions and citations, its own
citations, the documents, jurisdictions and document classes that ground it,
and the notes field. The model is asked to order term types by category and
then by the number of documents that ground them, and to aim for completeness.

Our first run cited section labels from memory and got them wrong in
9\% of citations, with 65 term types. We therefore append to the user prompt a
\emph{manifest}: one line per document listing its index, title and every exact
section label, with the instruction that no other labels exist. The manifest is
generated from the corpus at run time, so it cannot drift from the file the
model sees.

\subsection{One Call over a Cached Corpus}
\label{sec:method:run}

We place the whole corpus in the system prompt, after the stable instructions
and behind a prompt-cache breakpoint with a one-hour lifetime; the task prompt
is the user message. Every later task, and every prompt iteration, then pays
only for the small uncached part. This choice is also a methodological one.
Quality rule four asks the model to recognise that \src{MiCA}'s ``referencing
the value of one official currency'' and the \src{GENIUS Act}'s ``fixed amount
of monetary value'' answer the same question. A retriever would have to
anticipate that two passages that share few words define the same dimension;
a model that sees all 24 documents at once does not.

A single call to \texttt{claude-haiku-5-5} at high effort produced the output
we analyse. We disable adaptive thinking because on a first attempt the model
spent more than 100,000 thinking tokens on this corpus without emitting any
text, whereas without thinking it answers directly. We stream the response, allow up
to 120,000 output tokens, and parse the JSON; if a response is cut off, the
parser salvages every complete term type and flags the result as truncated.
The run reported here ended normally after 302 seconds with 85,729 output
tokens, 2,542 uncached input tokens for the task prompt and manifest, and
235,298 input tokens read from the cache.

\subsection{Citation Verification}
\label{sec:method:verify}

A quote that the model labels verbatim may still be invented, compressed, or
attached to the wrong section, so we check all of them. For each citation we
look up the document and section, normalise both the quote and the section text
to lower-case alphanumerics, and accept the quote if it is a substring. If it is
not, we split the quote into 12-character shingles and accept it when at least
80\% of them occur inside one window of the section twice the quote's length.
Localising the matches tolerates footnote markers, hyphenation and the spacing
artefacts of PDF extraction, and rejects a paraphrase stitched from phrases
that occur in different places. When a quote fails under the cited label but passes under another label of
the same document, we \emph{relabel} it, correcting the section, keeping the
model's label alongside, and counting the citation as verified. Everything else is rejected and listed in a report.
Every citation in the published output carries a \texttt{verified} flag, and
the rejection rate is our measure of desideratum~(D1), faithfulness.
Verification establishes that the quote exists where the model says it does;
it does not establish that the quote supports the dimension. We return to this
limit in \cref{sec:exp:rq1}.

\subsection{From Term Types to a Hierarchy}
\label{sec:method:hierarchy}
\todo{Task B: for each term type, a probability distribution over more general
term types in the same category (plus \emph{none}), with rationale and
citations, asked against the same cached corpus plus the Task A catalogue as a
second cached block. Describe once the Task B run is reviewed.}
